\chapter{Derivaciones y detalles matemáticos}
\label{app:derivaciones}

\section{Suavidad de una señal sobre grafo}
Sea $\mathcal{G}=(\mathcal{V},\mathcal{E},\mathbf{W})$ un grafo no dirigido con matriz de pesos simétrica. Para una señal $\mathbf{x}\in\mathbb{R}^{N}$, la energía de Dirichlet sobre el grafo se escribe como
\begin{equation}
\mathbf{x}^{\top}\mathbf{L}\mathbf{x}=\frac{1}{2}\sum_{i,j} W_{ij}(x_i-x_j)^2.
\end{equation}
Esta identidad muestra que la regularización laplaciana penaliza diferencias grandes entre nodos altamente conectados. En interpolación EEG, esta hipótesis equivale a suponer que electrodos cercanos o funcionalmente relacionados no deberían presentar cambios arbitrariamente abruptos en ausencia de evidencia medida.

\section{Interpolación con canales observados y faltantes}
Separando el conjunto de canales en observados $\mathcal{K}$ y faltantes $\mathcal{U}$, el problema instantáneo puede expresarse como
\begin{equation}
\hat{\mathbf{x}}=\arg\min_{\mathbf{z}}\|\mathbf{M}(\mathbf{z}-\mathbf{x})\|_2^2+\alpha \mathbf{z}^{\top}\mathbf{L}\mathbf{z},
\end{equation}
donde $\mathbf{M}$ es una máscara diagonal que conserva únicamente canales observados. Esta formulación permite combinar fidelidad a datos medidos con suavidad en grafo.

\section{Regularización temporal}
Para una matriz de señales $\mathbf{X}\in\mathbb{R}^{N\times T}$, TRSS introduce un término temporal. Una forma general es
\begin{equation}
\hat{\mathbf{X}}=\arg\min_{\mathbf{Z}}\|\mathbf{M}(\mathbf{Z}-\mathbf{X})\|_F^2+\alpha\operatorname{tr}(\mathbf{Z}^{\top}\mathbf{L}\mathbf{Z})+\beta\|\mathbf{Z}\mathbf{D}_t\|_F^2,
\end{equation}
donde $\mathbf{D}_t$ representa diferencias finitas en el eje temporal. El parámetro $\alpha$ controla suavidad espacial y $\beta$ controla continuidad temporal. Cuando $\beta=0$, el problema se reduce a regularización espacial independiente por muestra temporal.

\section{Interpretación espectral}
Si $\mathbf{L}=\mathbf{U}\mathbf{\Lambda}\mathbf{U}^{\top}$, entonces
\begin{equation}
\mathbf{x}^{\top}\mathbf{L}\mathbf{x}=\sum_{\ell=1}^{N}\lambda_{\ell}\hat{x}_{\ell}^{2}, \quad \hat{\mathbf{x}}=\mathbf{U}^{\top}\mathbf{x}.
\end{equation}
La regularización laplaciana penaliza componentes asociadas a autovalores grandes, análogas a frecuencias espaciales altas sobre el grafo. Esta lectura justifica entender Tikhonov y TRSS como filtros paso-bajo espaciales, complementados por un filtro temporal cuando se incorpora $\beta$.

\section{Criterio de estabilidad numérica}
El paso de gradiente debe respetar el Lipschitz de la función objetivo. En términos cualitativos, aumentar $\alpha$ o $\beta$ incrementa la curvatura del problema, por lo que se requiere un paso menor para garantizar convergencia. Este comportamiento motiva la búsqueda automatizada de hiperparámetros y el registro de configuraciones óptimas por escenario.
